\documentclass[a4paper,landscape,8pt]{extarticle}
\usepackage[landscape,margin=0.7cm,top=0.5cm,bottom=0.7cm]{geometry}
\usepackage[utf8]{inputenc}
\usepackage[T1]{fontenc}
\usepackage[ngerman]{babel}
\usepackage{helvet}
\renewcommand{\familydefault}{\sfdefault}
\usepackage{xcolor}
\usepackage{tikz}
\usetikzlibrary{positioning}
\usepackage{enumitem}
\usepackage{multicol}
\usepackage{array}
\usepackage{fancyhdr}
\usepackage{amssymb}

% Farben
\definecolor{familygreen}{HTML}{43A047}
\definecolor{appleblue}{HTML}{007AFF}
\definecolor{vhsorange}{HTML}{FF9800}
\definecolor{lightgray}{HTML}{F5F5F5}
\definecolor{bordergray}{HTML}{BBBBBB}

\pagestyle{fancy}
\fancyhf{}
\renewcommand{\headrulewidth}{0pt}
\fancyfoot[C]{\footnotesize Seite \thepage{} von 3}

\setlength{\parindent}{0pt}
\setlength{\parskip}{2pt}

% Kompakte Listen
\setlist{nosep, leftmargin=1em, itemsep=0pt, topsep=1pt}

\begin{document}

% ============================================================================
% SEITE 1: INFORMATIONSBLATT
% ============================================================================

\noindent{\Large\bfseries Hilfe holen: Wer kann mir helfen?} \hfill
{\footnotesize\textbf{Lernziele:} 3 Anlaufstellen | Problem beschreiben | Apple Support}

\vspace{1mm}
\hrule
\vspace{1mm}

% Drei-Spalten-Layout
\begin{minipage}[t]{0.32\linewidth}
\begin{center}
\colorbox{familygreen!15}{%
\begin{minipage}{0.95\linewidth}
\vspace{2mm}
\begin{center}
{\Large\bfseries\textcolor{familygreen}{1. Vertrauensperson}}\\[1pt]
{\footnotesize\itshape Familie \& Freunde}
\end{center}
\vspace{2mm}
\end{minipage}%
}
\end{center}

\vspace{1mm}

\textbf{Erste Anlaufstelle:}\\
Jemand aus Familie oder Freundeskreis, der sich auskennt.

\textbf{Vorteile:}
\begin{itemize}
\item Kennt dich und dein Gerät
\item Kann direkt helfen (vor Ort)
\item Kostenlos
\item Geduldig erklären
\end{itemize}

\textbf{Gut für:}
\begin{itemize}
\item Kleine Fragen (``Wo finde ich...?'')
\item Einstellungen ändern
\item Apps installieren
\item Allgemeine Tipps
\end{itemize}

\vspace{2mm}
\fcolorbox{familygreen}{familygreen!5}{%
\begin{minipage}{0.92\linewidth}
\footnotesize
\textbf{Tipp:} Leg dir eine feste Vertrauensperson zu -- frag \textbf{eine} Person regelmäßig, nicht jedes Mal jemand anderes!
\end{minipage}%
}

\vspace{2mm}

\textbf{Wichtig:}\\
Überfordere die Person nicht! Bei größeren Problemen lieber Apple Support.

\end{minipage}%
\hfill%
\begin{minipage}[t]{0.32\linewidth}
\begin{center}
\colorbox{appleblue!15}{%
\begin{minipage}{0.95\linewidth}
\vspace{2mm}
\begin{center}
{\Large\bfseries\textcolor{appleblue}{2. Apple Support}}\\[1pt]
{\footnotesize\itshape Der offizielle Weg}
\end{center}
\vspace{2mm}
\end{minipage}%
}
\end{center}

\vspace{1mm}

\textbf{Telefon (kostenlos):}\\
{\Large\bfseries 0800 6645 451}\\
{\footnotesize Mo--Fr 8--20 Uhr, Sa 9--18 Uhr}

\textbf{Weitere Wege:}
\begin{itemize}
\item \textbf{Website:} support.apple.com/de-de
\item \textbf{App:} ``Apple Support'' (App Store)
\item \textbf{Chat:} Über Website oder App
\item \textbf{Store:} Termin im Apple Store
\end{itemize}

\textbf{Gut für:}
\begin{itemize}
\item Technische Probleme
\item Garantiefälle
\item Wenn nichts mehr hilft
\item Sicherheitsfragen
\end{itemize}

\vspace{2mm}
\fcolorbox{appleblue}{appleblue!5}{%
\begin{minipage}{0.92\linewidth}
\footnotesize
\textbf{Kostenlos?} Beratung ja! Reparaturen können kosten. Frag vorher nach!
\end{minipage}%
}

\end{minipage}%
\hfill%
\begin{minipage}[t]{0.32\linewidth}
\begin{center}
\colorbox{vhsorange!15}{%
\begin{minipage}{0.95\linewidth}
\vspace{2mm}
\begin{center}
{\Large\bfseries\textcolor{vhsorange}{3. Kurse \& Angebote}}\\[1pt]
{\footnotesize\itshape Lernen in der Nähe}
\end{center}
\vspace{2mm}
\end{minipage}%
}
\end{center}

\vspace{1mm}

\textbf{Volkshochschule (VHS):}
\begin{itemize}
\item Kurse für Anfänger
\item Oft speziell für Senioren
\item Günstig (oft 30--60 Euro)
\item In deiner Stadt
\item Website: www.vhs.de
\end{itemize}

\textbf{Weitere Angebote:}
\begin{itemize}
\item \textbf{Seniorenbüro} der Stadt
\item \textbf{Mehrgenerationenhäuser}
\item \textbf{Bibliotheken} (oft kostenlos)
\item \textbf{Kirchengemeinden}
\end{itemize}

\textbf{Gut für:}
\begin{itemize}
\item Grundlagen lernen
\item Neue Funktionen entdecken
\item Andere Lernende treffen
\item Regelmäßig üben
\end{itemize}

\vspace{2mm}
\fcolorbox{vhsorange}{vhsorange!5}{%
\begin{minipage}{0.92\linewidth}
\footnotesize
\textbf{Tipp:} Such nach ``Smartphone-Kurs Senioren'' + dein Ort!
\end{minipage}%
}

\end{minipage}

\vspace{1mm}

% Zusammenfassung
\begin{center}
\fcolorbox{bordergray}{lightgray}{%
\begin{minipage}{0.98\linewidth}
\centering
\vspace{1.5mm}
\textbf{Die 3 Anlaufstellen:}
\textcolor{familygreen}{\textbf{1. Vertrauensperson}} (kleine Fragen) ~$|$~
\textcolor{appleblue}{\textbf{2. Apple Support}} (technische Probleme, Tel.: 0800 6645 451) ~$|$~
\textcolor{vhsorange}{\textbf{3. Kurse/VHS}} (Neues lernen)
\vspace{1.5mm}
\end{minipage}%
}
\end{center}

\newpage

% ============================================================================
% SEITE 2: ÜBUNGEN
% ============================================================================

\begin{center}
{\LARGE\bfseries Übungen: Hilfe holen}\\[3pt]
{\normalsize Schau dir Seite 1 an und beantworte die Fragen}
\end{center}

\vspace{3mm}
\hrule
\vspace{4mm}

\begin{multicols}{2}

\textbf{\large Teil A: Die 3 Anlaufstellen}

Richtig oder Falsch? Kreuze an.

\vspace{2mm}

\renewcommand{\arraystretch}{1.6}
\begin{tabular}{@{}p{8cm}cc@{}}
& R & F \\
1. Der Apple Support ist immer kostenpflichtig. & $\square$ & $\square$ \\
2. Die VHS bietet Kurse für Senioren an. & $\square$ & $\square$ \\
3. Die Vertrauensperson sollte jedes Mal wechseln. & $\square$ & $\square$ \\
4. Die Apple Support Telefonnummer ist kostenlos. & $\square$ & $\square$ \\
5. Bibliotheken bieten manchmal kostenlose Kurse. & $\square$ & $\square$ \\
\end{tabular}
\renewcommand{\arraystretch}{1}

\vspace{6mm}

\textbf{\large Teil B: Wer hilft wobei?}

Verbinde Problem und Anlaufstelle:

\vspace{2mm}

\renewcommand{\arraystretch}{1.8}
\begin{tabular}{@{}p{5cm}p{4cm}@{}}
``Wie installiere ich WhatsApp?'' & $\circ$ ~~~~ $\circ$ Apple Support \\
``Mein iPhone startet nicht mehr'' & $\circ$ ~~~~ $\circ$ Vertrauensperson \\
``Ich möchte mehr lernen'' & $\circ$ ~~~~ $\circ$ VHS-Kurs \\
``Garantiefall -- Display kaputt'' & $\circ$ ~~~~ $\circ$ Apple Store \\
\end{tabular}
\renewcommand{\arraystretch}{1}

\vspace{6mm}

\textbf{\large Teil C: Problem richtig beschreiben}

Wenn du Hilfe brauchst, solltest du sagen können:

\vspace{2mm}
\begin{enumerate}
\item \textbf{Was} geht nicht? \dotfill
\item \textbf{Seit wann}? \dotfill
\item \textbf{Was hast du versucht}? \dotfill
\item \textbf{Fehlermeldung}? \dotfill
\end{enumerate}

\columnbreak

\textbf{\large Teil D: Wichtige Nummern}

Fülle aus:

\vspace{3mm}

Apple Support Telefon: \dotfill

\vspace{2mm}

Öffnungszeiten: Mo--Fr \dotfill Uhr

\vspace{2mm}

VHS-Website: \dotfill

\vspace{6mm}

\textbf{\large Teil E: Mein Problem beschreiben}

Denke an ein echtes oder ausgedachtes Problem. Beschreibe es so, dass jemand dir helfen kann:

\vspace{3mm}

\textbf{Was geht nicht?}\\
\rule{\linewidth}{0.4pt}

\vspace{3mm}

\textbf{Seit wann?}\\
\rule{\linewidth}{0.4pt}

\vspace{3mm}

\textbf{Was hast du schon versucht?}\\
\rule{\linewidth}{0.4pt}

\vspace{3mm}

\textbf{Gibt es eine Fehlermeldung?}\\
\rule{\linewidth}{0.4pt}

\vspace{5mm}

\textbf{Bonus: Meine Vertrauensperson}

\vspace{2mm}

Name: \dotfill

\vspace{2mm}

Telefon: \dotfill

\end{multicols}

\newpage

% ============================================================================
% SEITE 3: CHECKLISTE
% ============================================================================

\begin{center}
{\LARGE\bfseries Checkliste: Das weiß ich jetzt}\\[3pt]
{\normalsize Geh die Liste durch und hake ab, was du sicher weißt}
\end{center}

\vspace{3mm}
\hrule
\vspace{6mm}

% Drei Boxen nebeneinander
\begin{minipage}[t]{0.31\linewidth}
\begin{center}
\fcolorbox{familygreen}{familygreen!10}{%
\begin{minipage}{0.92\linewidth}
\vspace{3mm}
\begin{center}
{\large\bfseries\textcolor{familygreen}{Vertrauensperson}}
\end{center}
\vspace{2mm}
\begin{itemize}[label=$\square$, itemsep=5pt]
\item Ich habe eine \textbf{feste Person}
\item Ich weiß, wofür ich sie \textbf{frage}
\item Ich \textbf{überfordere} sie nicht
\item Für größere Probleme: Apple
\end{itemize}
\vspace{3mm}
\end{minipage}%
}
\end{center}
\end{minipage}%
\hfill%
\begin{minipage}[t]{0.31\linewidth}
\begin{center}
\fcolorbox{appleblue}{appleblue!10}{%
\begin{minipage}{0.92\linewidth}
\vspace{3mm}
\begin{center}
{\large\bfseries\textcolor{appleblue}{Apple Support}}
\end{center}
\vspace{2mm}
\begin{itemize}[label=$\square$, itemsep=5pt]
\item Ich kenne die \textbf{Telefonnummer}
\item Ich weiß: \textbf{Beratung kostenlos}
\item Ich kann mein Problem \textbf{beschreiben}
\item Ich kenne die \textbf{4 Kontaktwege}
\end{itemize}
\vspace{3mm}
\end{minipage}%
}
\end{center}
\end{minipage}%
\hfill%
\begin{minipage}[t]{0.31\linewidth}
\begin{center}
\fcolorbox{vhsorange}{vhsorange!10}{%
\begin{minipage}{0.92\linewidth}
\vspace{3mm}
\begin{center}
{\large\bfseries\textcolor{vhsorange}{Kurse \& Angebote}}
\end{center}
\vspace{2mm}
\begin{itemize}[label=$\square$, itemsep=5pt]
\item Ich kenne die \textbf{VHS}
\item Ich weiß von \textbf{Seniorenbüros}
\item Bibliotheken haben \textbf{Kurse}
\item Ich kann nach Kursen \textbf{suchen}
\end{itemize}
\vspace{3mm}
\end{minipage}%
}
\end{center}
\end{minipage}

\vspace{8mm}

% Gesamtverständnis
\begin{center}
\fcolorbox{black}{lightgray}{%
\begin{minipage}{0.95\linewidth}
\vspace{4mm}
\begin{center}
{\large\bfseries Problem richtig beschreiben}
\end{center}
\vspace{3mm}
\begin{multicols}{2}
\begin{itemize}[label=$\square$, itemsep=6pt]
\item Ich sage, \textbf{was} nicht geht
\item Ich sage, \textbf{seit wann}
\item Ich sage, was ich \textbf{versucht} habe
\item Ich notiere \textbf{Fehlermeldungen}
\item Ich \textbf{traue mich}, um Hilfe zu bitten
\end{itemize}
\end{multicols}
\vspace{3mm}
\end{minipage}%
}
\end{center}

\vspace{8mm}

% Notfall-Notizen
\begin{center}
\fcolorbox{bordergray}{white}{%
\begin{minipage}{0.95\linewidth}
\vspace{4mm}
\begin{center}
{\large\bfseries Meine Hilfe-Kontakte}\\[3pt]
{\footnotesize (Trage hier deine wichtigsten Anlaufstellen ein)}
\end{center}
\vspace{4mm}
\renewcommand{\arraystretch}{2}
\begin{tabular}{@{}p{0.45\linewidth}p{0.45\linewidth}@{}}
\textbf{Meine Vertrauensperson:} & \textbf{Apple Support:} \\
Name: \dotfill & Tel.: 0800 6645 451 \\[2mm]
Tel.: \dotfill & Website: support.apple.com/de-de \\[2mm]
\textbf{VHS in meiner Stadt:} & \textbf{Sonstiges Angebot:} \\
\dotfill & \dotfill \\[2mm]
Tel./Web: \dotfill & Tel./Web: \dotfill \\
\end{tabular}
\renewcommand{\arraystretch}{1}
\vspace{4mm}
\end{minipage}%
}
\end{center}

\vfill

\begin{center}
\footnotesize\textit{Stand: \today{} -- Es ist keine Schwäche, um Hilfe zu bitten! Die richtige Anlaufstelle spart Zeit und Nerven.}
\end{center}

\end{document}
